\scanpage{38}
\begin{center}6 - 27 -\hfill n\textsuperscript{o} 392\end{center}

12) Démontrer la prop. 5 du n\textsuperscript{o} 8 en utilisant le chap. III, \S\ 2, n\textsuperscript{o} 8, cor. 3 du th. 1.

13) Etendre le cor. de la prop. 7 du n\textsuperscript{o} 8 au cas où l'anneau local séparé et complet $A$ n'est pas intègre, (utiliser \uline{Alg.}, chap. VIII, \S\ 6, exerc. 6 b)).

14) Soient $A$ un anneau local complet noethérien, dont le corps résiduel est de caractéristique $p>0$. Dans l'anneau de séries formelles $A[[T]]$, on considère les éléments $\omega_n=(1-T)^{p^n}$ et $\gamma_n=1-\omega_n$ pour tout entier $n>0$. Montrer que $\gamma_n$ est, au signe près, un polynôme distingué, (n\textsuperscript{o} 8) ; en déduire que $A_n=A[[T]]/(\gamma_n)$ s'identifie à l'algèbre sur $A$ du groupe $G_n=Z/p^nZ$. Montrer que l'intersection des idéaux principaux $(\gamma_n)$ est réduite à $0$ ; en déduire que $A[[T]]$ s'identifie à la limite projective $\varprojlim A_n$.

15) Soient $K$ un corps complet pour une valuation discrète $v$, $A$ l'anneau de la valuation, $k$ son corps résiduel, $P$ un polynôme de $A[X_1,\ldots,X_n]$ de degré total $d$ ayant la propriété suivante : il existe une extension algébrique $K'$ de $K$ telle que dans $K'[X_1,\ldots,X_n]$, $P$ soit produit de polynômes de degré total $1$. On suppose en outre qu'il existe deux polynômes $Q$, $R$ de $A[X_1,\ldots,X_n]$ tels que $Q$ soit de degré total $s$ et contienne un monôme $aX_1^s$ où $\varphi(a)\ne 0$, ($\varphi$ désignant l'homomorphisme canonique $A\longrightarrow k$), que $R$ soit de degré $\leq d-s$, et que l'on ait $\overline{P}=\overline{Q}.\overline{R}$, (notations du chap. III, \S\ 4). Montrer qu'il existe alors dans $A[X_1,\ldots,X_n]$ deux polynômes $Q_0$, $R_0$ de degrés respectifs $s$, $d-s$, tels que $P=Q_0R_0$, $\overline{Q}=\overline{Q}_0$, $\overline{R}=\overline{R}_0$ et que $Q_0$ contienne un monôme $a_0X_1^s$ avec $\varphi(a)=\varphi(a_0)$.

(Considérer $P$, $Q$, $R$ comme des polynômes en $X_1$ à coefficients dans l'anneau $B$, complété de $A[X_2,\ldots,X_n]$ pour la valuation obtenue en prolongeant $v$ suivant la méthode du chap. VI, \S\ 10, n\textsuperscript{o}\quad ; appliquer ensuite le lemme de Hensel ; utiliser enfin l'hypothèse initiale sur $P$).

16) Soient $B$ un anneau de valuation discrète, dont le corps résiduel $k$ est fini et n'est pas un corps premier ; soient $k_0$ le sous-corps premier de $k$, et soit $A$ le sous-anneau de $B$ formé des éléments dont la classe dans le corps résiduel appartient à $k_0$. Soient $\pi$ une uniformisante de $B$, et $(\unsure{\vartheta|\theta}_i)_{1\leq i\leq m}$ un système d'éléments inversibles de $B$ tels que les classes $\overline{\unsure{\vartheta|\theta}}_i$ mod. $\pi$ des $\unsure{\vartheta|\theta}_i$

\margin{37}

